
\documentclass[border=4pt]{standalone}
\usepackage{polyglossia}
\setdefaultlanguage{russian}
\setotherlanguage{english}
\usepackage{fontspec}
\IfFontExistsTF{CMU Serif}{\setmainfont{CMU Serif}}{\setmainfont{Liberation Serif}}
\IfFontExistsTF{CMU Sans Serif}{\setsansfont{CMU Sans Serif}}{\setsansfont{Liberation Sans}}
\IfFontExistsTF{CMU Typewriter Text}{\setmonofont{CMU Typewriter Text}}{\setmonofont{DejaVu Sans Mono}}
\usepackage{amsmath,amssymb,mathtools,bm}
\usepackage{xcolor}
\definecolor{accent}{HTML}{1F4E79}
\definecolor{accent2}{HTML}{B85450}
\definecolor{shade}{HTML}{F4F1EA}
\colorlet{prosvBlue}{accent}
\colorlet{prosvRed}{accent2}
\colorlet{prosvLight}{accent!12}
\colorlet{prosvCream}{shade}
\definecolor{prosvGold}{HTML}{E8B647}
\definecolor{prosvGreen}{HTML}{6A9F58}
\colorlet{prosvGray}{black!55}
\usepackage{tikz}
\usetikzlibrary{calc, arrows.meta, decorations.pathreplacing,
                positioning, shapes.geometric, shapes.misc, fit,
                backgrounds}
\usepackage{pgfplots}
\pgfplotsset{compat=1.18}
\newcommand{\R}{\mathbb{R}}
\newcommand{\N}{\mathbb{N}}
\newcommand{\Z}{\mathbb{Z}}
\newcommand{\eps}{\varepsilon}
\newcommand{\norm}[1]{\left\|#1\right\|}
\newcommand{\abs}[1]{\left|#1\right|}
\newcommand{\NOD}{\gcd}
\newcommand{\Eul}{\varphi}
\newcommand{\Zn}[1]{\mathbb{Z}_{#1}}
\newcommand{\modn}[1]{\,(\mathrm{mod}\,#1)}
\newcommand{\term}[1]{\textcolor{accent2}{\textbf{#1}}}
\newcommand{\engterm}[1]{\textit{#1}}
\DeclareMathOperator*{\argmin}{arg\,min}
\DeclareMathOperator*{\argmax}{arg\,max}
\begin{document}

\begin{tikzpicture}
\begin{axis}[
  width=11cm, height=6cm,
  xlabel={$N$}, ylabel={число простых $\le N$},
  xmin=1, xmax=1000, ymin=0, ymax=200,
  axis lines=left,
  grid=major, grid style={prosvGray!30},
  legend pos=north west, legend cell align=left,
  xlabel style={font=\small}, ylabel style={font=\small},
  tick label style={font=\scriptsize},
  every axis plot/.append style={very thick}
]

\addplot[prosvRed, domain=2:1000, samples=80] {x/ln(x)};
\addlegendentry{$\dfrac{N}{\ln N}$ (приближение)};

\addplot+[
  only marks, mark=*, prosvBlue, mark size=1.4pt
] coordinates {
  (10,4) (50,15) (100,25) (200,46) (300,62) (400,78) (500,95)
  (600,109) (700,125) (800,139) (900,154) (1000,168)
};
\addlegendentry{точное значение $\pi(N)$};
\end{axis}
\end{tikzpicture}
\end{document}
